\newpage
%% [migrated] heading absorbed into chapter wrapper: Strong Interactions beyond the SM



\begin{itemize}
    \item While the main focus of LFT calculations has been on QCD, the same (or similar) techniques are useful for studying strongly-interacting field theories more generally both in particle/nuclear physics and more generally
    \item One class of such investigations is in near-conformal gauge theories that may play a role in BSM scenarios for explaining the Higgs boson as a composite bound state of an additional gauge theory with dynamics at a high scale - technicolour
    \begin{itemize}
        \item The simplest implementations of technicolour are in conflict with precision electroweak sector tests, however in ``walking technicolour'' theories (see Fig.~\ref{fig:walking}), the content of the theory is set up such that the coupling becomes almost scale independent for a range of scales (this separation of scales allows for EW constraints to be avoided)
        \item Perturbatively, the YM $\beta$ function has an IR fixed point depending on number of colours and fermion flavours ($N_f$ fundamental rep assumed, see Fig.~\ref{fig:conformal_fp})
\begin{figure}
\begin{subfigure}{0.4\textwidth}
\redrawcompare{fig-notes-042.jpg}{fig-redraw-042.pdf}
\end{subfigure}
\begin{subfigure}{0.4\textwidth}
\redrawcompare{fig-notes-040.jpg}{fig-redraw-040.pdf}
\end{subfigure}
\caption{The behavior of the coupling and beta function as a function of scale for walking technicolour.
\label{fig:walking}}
\end{figure}
\begin{equation}
    \beta(g) = -\beta_0 g^3 - \beta_1 g^5 + \ldots,
\end{equation}
with $\beta_0 = \frac{11}{3}C_r - \frac{4}{3} T_r N_f$ and $\beta_1 = \frac{34}{3}C_r^2 - \frac{20}{3} C_r T_r N_f - 4 C_r T_r N_f$. However, the coupling $g_*$ at the fixed point may be outside the regime of perturbative control, so conclusions are ambiguous
    \end{itemize}

\begin{figure}
\redrawcompare{fig-notes-041.jpg}{fig-redraw-041.pdf}
\caption{Various trajectories for the coupling $g$ for different values of the number of colours and fermion flavours, showing the possibility of an IR fixed point.
\label{fig:conformal_fp}}
\end{figure}

    \item Additionally an actual perturbative fixed point would not be consistent with phenomenology of confinement so what is being sought is a theory that comes very close to having a fixed point but then runs away from it
    \item Requires non-perturbative study of the scale dependence of the coupling using LQFT methods. This is complicated because for phenomenology, an exponential separation of scales between $\Lambda_*$ and $\Lambda_{\text{IR}}$ is desired: difficult with hypercubic lattice
    \item A way around this issue is to use step-scaling to define the coupling in a twisted Schr\"odinger functional scheme (see Section \ref{sec:schrodinger_functional}).
    

    \item Studies have been performed in many different theories (e.g. $SU(N_c=2)$ with $N_f = 2$ adjoint fermion species) with some indications of interesting near-conformal behaviour. Results are clouded by potential for non-universal behaviour away from UV Gaussian fixed point.
    \item As well as seeing walking coupling, a theory for composite Higgs would need to produce a ``light'' scalar (since we now know the Higgs at the LHC is not pseudoscalar), so need to study the spectrum and interactions of potential BSM theories too
    \item Strongly-interacting theories may also be the underlying theory of self-interacting dark matter
    \begin{itemize}
        \item In this context, we want to understand the space of behaviours that occur in such theories. Are features of QCD generic? E.g. nuclei? Are there features that occur that are not just modified versions of QCD?
        \item Quite a few studies of $SU(N_c = 2)$ and $SU(N_c = 4)$ as potential models of dark matter (or some component thereof) either with or without fermions
    \end{itemize}
    \item LQFT technology developed for QCD such as HMC has been used in calculations of e.g. graphene and nonrelativistic fermion systems such as nuclear lattice effective theory (NLEFT)
\end{itemize}
